% =========================================================
% 3.13 Adder
% =========================================================

\section{অ্যাডার (\textenglish{Adder})}

কম্পিউটার বিভিন্ন ধরনের গাণিতিক কাজ করলেও তার অভ্যন্তরে সব তথ্য
\(0\) ও \(1\)-এর সমন্বয়ে প্রক্রিয়াকরণ করা হয়। তাই আমরা কম্পিউটার বা
ক্যালকুলেটরে দশমিক সংখ্যা লিখলেও ডিজিটাল সার্কিট সরাসরি দশমিক অঙ্ক নিয়ে
যোগ করে না; সংখ্যাগুলো আগে উপযুক্ত
\textenglish{binary representation}-এ প্রকাশিত হয় এবং এরপর
\textenglish{binary addition} সম্পন্ন হয়।

ধরা যাক, আমরা লিখলাম—

\[
4+7
\]

এখানে,

\[
(4)_{10}=(0100)_2,
\qquad
(7)_{10}=(0111)_2
\]

ডিজিটাল সার্কিটের ভেতরে মূলত যোগ হয়—

\[
\begin{array}{r}
0100\\[-1mm]
+\,0111\\
\hline
1011
\end{array}
\]

এবং,

\[
(1011)_2=(11)_{10}
\]

অতএব ব্যবহারকারী ফলাফল হিসেবে \(11\) দেখে।

\par\begingroup

\textbf{মনে রাখবে:}
\textenglish{Adder} সরাসরি দশমিক অঙ্কের অর্থ বোঝে না।
ডিজিটাল সিস্টেমে সংখ্যা আগে \textenglish{binary/code} আকারে প্রকাশিত হয়;
\textenglish{Adder} সেই বাইনারি বিটগুলোর উপর যোগ সম্পন্ন করে।
\par\endgroup



% =========================================================
% 3.13.1 What is Adder?
% =========================================================

\subsection{অ্যাডার কী?
(\textenglish{What is Adder?})}

\par\begingroup

যে সমবায় লজিক সার্কিট
(\textenglish{Combinational Logic Circuit})
বাইনারি বিটসমূহের যোগ সম্পন্ন করে এবং প্রয়োজন অনুযায়ী
\textenglish{Sum} ও \textenglish{Carry} তৈরি করে, তাকে
\textbf{অ্যাডার (\textenglish{Adder})} বলে।
\par\endgroup


\textenglish{Adder} হলো ডিজিটাল কম্পিউটারের
\textenglish{arithmetic operation}-এর একটি মৌলিক সার্কিট।
\textenglish{Processor}-এর
\textenglish{ALU (Arithmetic Logic Unit)}-তে বাইনারি যোগ সম্পন্ন করার
জন্য \textenglish{Adder} ব্যবহৃত হয়।

\begin{center}
\begin{chThreeDiagram}
\begin{tikzpicture}[
    >=Stealth,
    line width=0.8pt,
    node distance=1.0cm
]

\node[
    draw,
    rounded corners=2pt,
    fill=ICTLightBlue,
    minimum width=3.1cm,
    minimum height=0.9cm
] (in) {\textenglish{Binary Inputs}};

\node[
    draw,
    rounded corners=2pt,
    fill=ICTLightGreen,
    minimum width=2.5cm,
    minimum height=0.9cm,
    right=of in
] (adder) {\textbf{\textenglish{Adder}}};

\node[
    draw,
    rounded corners=2pt,
    fill=ICTLightOrange,
    minimum width=3.0cm,
    minimum height=0.9cm,
    right=of adder
] (out) {\textenglish{Sum + Carry}};

\draw[->] (in) -- (adder);
\draw[->] (adder) -- (out);

\end{tikzpicture}
\end{chThreeDiagram}
\end{center}


% =========================================================
% 3.13.2 Why Adder is Needed
% =========================================================

\subsection{কেন অ্যাডার প্রয়োজন?}

কম্পিউটারের বহু \textenglish{arithmetic operation}-এর সঙ্গে যোগের প্রত্যক্ষ
বা পরোক্ষ সম্পর্ক রয়েছে। তাই \textenglish{Adder} শুধু \(+\) চিহ্নের যোগের
জন্যই নয়, বিভিন্ন \textenglish{digital arithmetic circuit}-এর ভিত্তি
হিসেবেও ব্যবহৃত হয়।

\begin{itemize}[label={{\bulletfont\textbullet}}, leftmargin=*, nosep]

\item দুটি বা একাধিক বাইনারি বিট বা সংখ্যা যোগ করতে।

\item \textenglish{ALU}-তে
\textenglish{arithmetic operation} সম্পন্ন করতে।

\item বহু-বিট \textenglish{binary adder} তৈরিতে।

\item \textenglish{Address calculation} ও বিভিন্ন
\textenglish{digital arithmetic circuit}-এ।

\item উপযুক্ত পদ্ধতিতে \textenglish{subtraction}-সহ অন্যান্য
\textenglish{arithmetic operation} বাস্তবায়নে।

\end{itemize}


% =========================================================
% 3.13.3 Binary Addition and Carry
% =========================================================

\subsection{বাইনারি যোগ ও ক্যারি
(\textenglish{Binary Addition and Carry})}

\textenglish{Adder} বোঝার আগে বাইনারি যোগের চারটি মৌলিক নিয়ম জানা প্রয়োজন।

\begin{center}
\begin{tabular}{|c|c|c|c|}
\hline

\rowcolor{ICTLightBlue}

\(\mathbf A\)
&
\(\mathbf B\)
&
\textbf{\textenglish{Sum}}
&
\textbf{\textenglish{Carry}}
\\

\hline
0&0&0&0\\
\hline

0&1&1&0\\
\hline

1&0&1&0\\
\hline

1&1&0&1\\
\hline

\end{tabular}
\end{center}

প্রথম তিনটি ক্ষেত্রে কোনো \textenglish{Carry} তৈরি হয় না। কিন্তু—

\[
1+1=(10)_2
\]

এখানে ডান পাশের \(0\) হলো \textenglish{Sum}
এবং বাম পাশের \(1\) হলো \textenglish{Carry}।

অর্থাৎ,

\[
S=0,
\qquad
C=1
\]


% =========================================================
% 3.13.4 Types of Adder
% =========================================================

\subsection{অ্যাডারের প্রকারভেদ}

প্রাথমিকভাবে \textenglish{Adder} দুই প্রকার—

\begin{enumerate}[leftmargin=*,itemsep=2pt]

\item
\textbf{অর্ধযোগক (\textenglish{Half Adder})}

\item
\textbf{পূর্ণযোগক (\textenglish{Full Adder})}

\end{enumerate}

\textenglish{Half Adder} দুটি \textenglish{input bit} যোগ করতে পারে,
কিন্তু পূর্ববর্তী ধাপ থেকে আসা \textenglish{Carry} গ্রহণ করতে পারে না।

অন্যদিকে \textenglish{Full Adder} দুটি \textenglish{input bit}-এর সঙ্গে
পূর্ববর্তী ধাপ থেকে আসা \textenglish{Carry-in}-ও যোগ করতে পারে।


% =========================================================
% 3.13.5 Half Adder
% =========================================================

\subsection{অর্ধযোগক
(\textenglish{Half Adder})}

\par\begingroup

যে সমবায় লজিক সার্কিট দুটি ১-বিট বাইনারি
\textenglish{input} যোগ করে একটি \textenglish{Sum}
এবং একটি \textenglish{Carry output} প্রদান করে,
তাকে \textbf{অর্ধযোগক (\textenglish{Half Adder})} বলে।
\par\endgroup


\textenglish{Half Adder}-এর দুটি \textenglish{input} হলো

\[
A,\;B
\]

এবং দুটি \textenglish{output} হলো

\[
S,\;C
\]

যেখানে,

\[
S=\text{\textenglish{Sum}},
\qquad
C=\text{\textenglish{Carry}}
\]


% =========================================================
% Sum Boolean Expression
% =========================================================

\subsubsection*{\textenglish{Sum} \(S\)-এর বুলিয়ান রাশি নির্ণয়}

\textenglish{Truth Table}-এ \(S=1\) হয় যখন—

\[
A=0,\;B=1
\]

অথবা,

\[
A=1,\;B=0
\]

তাই,

\[
S=\bar A B+A\bar B
\]

আবার \textenglish{XOR}-এর সংজ্ঞা অনুযায়ী,

\[
A\oplus B
=
\bar A B+A\bar B
\]

অতএব,

\[
\boxed{S=A\oplus B}
\]


% =========================================================
% Carry Boolean Expression
% =========================================================

\subsubsection*{\textenglish{Carry} \(C\)-এর বুলিয়ান রাশি নির্ণয়}

\textenglish{Truth Table}-এ \(C=1\) হয় শুধু যখন—

\[
A=1,\;B=1
\]

তাই,

\[
C=A\cdot B
\]

অতএব,

\[
\boxed{C=AB}
\]


% =========================================================
% Half Adder Final 3-Column Figure
% =========================================================

\begin{center}

\includegraphics[
    width=0.98\textwidth,
    keepaspectratio
]{half_adder_box.png}

\end{center}

\vspace{-1mm}

\begin{center}
{\small
\textbf{চিত্র:}
\textenglish{Half Adder}-এর সত্যক সারণি এবং
যৌগিক ও মৌলিক গেট দিয়ে বাস্তবায়ন
}
\end{center}


% =========================================================
% Half Adder Limitation
% =========================================================

\par\begingroup


\textbf{হাফ অ্যাডারের সীমাবদ্ধতা:}

\textenglish{Half Adder} পূর্ববর্তী
\textenglish{bit-position} থেকে আসা
\textenglish{Carry-in} গ্রহণ করতে পারে না।

তাই \textenglish{Carry}-সহ তিনটি ১-বিট
\textenglish{input} যোগ করতে
\textenglish{Full Adder} প্রয়োজন।

\par\endgroup
